\section{Study 4: does benign consolidation actually produce this?}\label{sec:drift}

Studies 1--3 install the corruption by hand, which buys clean ground truth. It
also assumes something the literature repeats and, as far as we can tell, has
not measured: that summarisation drops caveats. We tested it, expecting
confirmation, and did not get it.

We ran a real consolidation chain over the same six source records, with no
instruction to preserve or drop anything and no mention of caveats, qualifiers,
or what matters: an analyst writes session notes; the notes are consolidated
into durable lessons; the durable memory is then re-summarised four times as it
ages and later cycles compete for context, with targets stepping
$130 \rightarrow 110 \rightarrow 90 \rightarrow 70$ characters. \FiveChains{}
chains, six generations each, across six models.

Survival is scored by a model that sees the original source record and all six
generations at once and reports, for each, whether the source's specific
constraint is still stated. Two of the six records carry no qualifier at all and
act as a false-positive check: the scorer identified them as unqualified in
every case and reported a constraint in $\FiveFalsePos$ judgments where none
exists.

\paragraph{The premise fails for negative outcomes.} A quantified negative
result is close to the last thing compression discards, though not immune. At
the tightest compression --- lessons cut to around 70 characters --- the
retention harm is still stated in \FiveNegResDPct\% of chains, ranging from
\FiveNegPerModelLo{} to \FiveNegPerModelHi{} across models. Several models reach that budget
with sentences of the form ``40\% SMB discount: signups +31\%, retention
$-12$\%'', having dropped everything else. The number is load-bearing enough to
be kept.

\paragraph{What does erode is scope, and normative force.} Scope restrictions go
roughly twice as fast: \FiveScopeResDPct\% survive at the same compression
against \FiveNegResDPct\% for the negative outcome. And within the corrupted
memory itself, separating the finding from the instruction shows the split
directly: at the consolidation stage \FiveNumCons{} chains state the negative
outcome and \FiveProhCons{} still state a prohibition; after four
re-summarisations, \FiveNumResD{} still state the outcome and only
\FiveProhResD{} state a prohibition. What consolidation removes is
not the fact that discounting hurt retention. It is the sentence saying not to
do it anyway.

\paragraph{Consequences for the threat model, including ours.} The corruption
Studies 1--3 use --- deleting an entire negative caveat --- is not what this
process produces. A benign-compression threat model should be built on the
operators compression actually applies: loss of scope, and loss of normative
constraint while the supporting number survives. Those are more dangerous than
they sound, because a lesson that retains its numbers looks well-evidenced. Our
own Study 3 says such a record draws less verification, not more.

We report this as a limitation on Studies 1--3 rather than repairing it here.
The allocation results stand: they concern which beliefs an agent checks given a
corrupted one, and hold for any corruption with the same decision geometry. What
does not stand is the claim that this particular corruption arises on its own.

\paragraph{A pre-registered rule that failed.} We registered that two scorers ---
a keyword rule and the model scorer --- must agree. That rule is mis-specified
and we report it rather than its output. Disagreements occur precisely where one
scorer sees a lost qualifier, so conditioning on agreement censors informatively
and returns 100\% survival at every generation by construction. The keyword rule
also turned out to have very low specificity on scope restrictions, firing on any
mention of the segment term, so it functions only as an upper bound. We report
the model scorer, validated against the unqualified controls, as primary, and
state the bound rather than presenting the agreeing subset as an estimate.
